\section{Archive layout and reproducibility}
\label{app:reproducibility}

The archive is a reviewable continuation against the pinned source,
not a remote repository update. It contains the complete article and
its figures, the modified \texttt{fast} source tree, an executable
baseline snapshot for comparisons, the required reference modules,
raw experimental evidence, and an integration patch. The existing
project license is retained with the source.

\subsection{File map}

\begin{description}[style=nextline,leftmargin=0pt,labelsep=0pt,labelwidth=0pt]
\item[\path{article/}]
  The main source \path{compiled_certificates.tex}, all included
  sections, the rendered PDF, generated tables, figure PDFs and PNGs,
  and \path{generate_artifacts.py}. No network access is required to
  rebuild the manuscript from the included sources.
\item[\path{fast/fastunknot/}]
  The full current runtime snapshot. The new modules are
  \path{raw_saturation.py} and \path{orbit_index.py}. Four existing
  files gain the optional group hook: \path{compressed_search.py},
  \path{group_certificate.py}, \path{recognize.py}, and
  \path{__main__.py}.
\item[\path{fast/tests/}]
  The maintained suite, including the three new test files
  \path{test_raw_saturation.py},
  \path{test_raw_saturation_integration.py}, and
  \path{test_orbit_index.py}.
\item[\path{fast/raw_saturation_research/}]
  The independent matrix/signed-word audit, four-mode benchmark,
  summary generator, and frozen initial-candidate snapshot. The
  initial snapshot explains the failed-probe optimization; its
  measurements are explicitly separated from the final candidate.
\item[\path{fast/orbit_index_research/}]
  The graph audit, complete membership benchmark, prepared rank/select
  benchmark, API documentation, and mathematical research notes.
\item[\path{fast/results/}]
  Exact inputs, raw timing samples, all outcomes, checked source hashes,
  logs, per-arm journals, and generated summaries. Only the completed
  final records supply headline values.
\item[\path{baseline/fast/} and \path{baseline/reports/}]
  The pinned runtime/test snapshot and its required reference modules.
  These support the paired group comparisons and the baseline suite
  without downloading the upstream repository.
\item[\path{reports/}]
  The exact reference Python modules from reports 24, 26, 28, and 45
  needed by existing tests. They are included because the maintained
  suite imports them from outside \texttt{fast}.
\item[\path{provenance/}]
  The Git-blob source manifest, archive-content hashes, environment
  record, baseline/final suite logs, and patch verification results.
\item[\path{integration.patch}]
  A source patch against the pinned repository paths, adding the two
  runtime modules, the optional hook, tests, and research drivers.
  The article and bulky experimental evidence are separately organized
  for repository placement.
\end{description}

The source manifest records the retrieval audit, including read-only
synthesis sources used during research. Those synthesis sources are not
all duplicated in the delivery. The archive-content manifest separately
identifies the files actually packaged. Baseline source bytes were checked
against their Git blob identifiers; current experimental source hashes
were checked before and after the measurements. The group source-hash
set excludes the unrelated orbit module, whose own evidence carries
complete separate hashes.

\subsection{Rebuilding and running the checks}

The runtime declares Python 3.10 or later and no mandatory package
dependencies. The recorded environment is Python 3.12.14, Regina 7.4.1,
NumPy 2.3.5, and Matplotlib 3.10.8. The full historical suite has tests
for optional native normal-surface functionality, so reproducing every
recorded test requires the supplied environment or the documented
Regina extra. The new saturation and orbit-index tests themselves use
the Python implementation and independent finite oracles.

From the archive's top directory, the principal commands are:
\begin{verbatim}
cd fast
python3 -B -m unittest discover -s tests -q
python3 -B raw_saturation_research/audit_epoch.py
python3 -B raw_saturation_research/benchmark.py \
  --mode audit --candidate-label final-preindex-guard \
  --output results/raw_saturation_audit.json
python3 -B raw_saturation_research/benchmark.py \
  --mode kernels --candidate-label final-preindex-guard \
  --output results/raw_saturation_kernels.json
python3 -B raw_saturation_research/benchmark.py \
  --mode stage --candidate-label final-preindex-guard \
  --output results/raw_saturation_stage.json
python3 -B raw_saturation_research/benchmark.py \
  --mode pipeline --candidate-label final-preindex-guard \
  --output results/raw_saturation_pipeline.json
python3 -B raw_saturation_research/summarize.py
python3 -B orbit_index_research/audit.py
python3 -B orbit_index_research/benchmark.py --rounds 7
python3 -B orbit_index_research/rank_select.py --rounds 7
\end{verbatim}
The group benchmark's default baseline location is
\path{../baseline/fast}. A separate checkout can be supplied with
\texttt{--baseline}. These commands replace the correspondingly named
result files; preserve the delivered records in another directory before
an independent replication. Run timing experiments without simultaneous
CPU-intensive tests. Random call order and controls reduce some timer
effects but do not turn a shared machine into an isolated laboratory.

From the archive's top directory, regenerate the numerical artifacts and PDF:
\begin{verbatim}
cd article
python3 generate_artifacts.py --results ../fast/results
latexmk -pdf -interaction=nonstopmode -halt-on-error \
  compiled_certificates.tex
\end{verbatim}
Two or more ordinary \texttt{pdflatex} passes can replace
\texttt{latexmk}. The preamble uses common TeX Live packages, with no
shell-escape requirement or external bibliography processor. The provided
PDF was rendered to page images and inspected for clipping, table
overflow, unresolved references, and figure readability.

\subsection{Review and integration order}

First review the mathematical theorem statements together with their
counterexamples. Next review the two runtime modules and the four small
option-plumbing changes. Apply the source patch to a clean checkout of
the pinned commit, or reconcile it manually with a newer upstream
revision; test commands should run from the \texttt{fast} directory.
The patch does not assign a new report number or overwrite the existing
synthesis manuscript. The standalone article can be moved to the
maintainer's chosen research location with all its included files.

The raw producer requires no certificate-verifier change. Its final
source proofs were checked by the old and current replayers. The orbit
index's documented factory independently verifies the entire supplied
trace, but its query arithmetic is an additional trusted implementation.
Its immutability contract concerns factory-created objects; manually
constructing private Python dataclass fields is outside that interface.
These are distinct trust boundaries and should remain distinct during
integration.

No general quasi-polynomial flag is enabled. The existing recognizer's
worst-case statement remains appropriate. The finite audit supplies no
claim that all possible mutants, diagrams, or interval systems have been
tested. The proofs establish the local statements; the recorded evidence
checks their implementation and the observed performance on specified
inputs.

\section{Complete group benchmark tables}
\label{app:full-tables}

In the following tables, U denotes \texttt{UNKNOT}, K denotes
\texttt{KNOTTED}, I denotes a positive-only group-stage
\texttt{INCONCLUSIVE}, and ? denotes full-pipeline \texttt{UNKNOWN}.
Times are milliseconds. A dash in a completion-time or ratio column
indicates an unavailable completed comparison; it is not zero. Observed
elapsed-time columns are populated only for unresolved arms.
Exact observed times and the reason for each unfinished call are retained
in the JSON and per-arm journals. All conclusive paired statuses agree.

\subsection{Group stage: reconstruction, search, and independent replay}

% Generated by generate_artifacts.py. This is a longtable, not a float.
\begingroup
\small
\setlength{\tabcolsep}{3pt}
\setlength{\LTcapwidth}{\linewidth}
\renewcommand{\arraystretch}{1.08}
\begin{longtable}{@{}>{\raggedright\arraybackslash}p{3.45cm}rrrrrrc@{}}
\caption{Full group-stage benchmark: all 28 inputs.}\label{tab:raw-stage-full}\\
\toprule
Input & $n$ & \shortstack{Old\\complete} & \shortstack{New\\complete} & \shortstack{Old\\observed} & \shortstack{New\\observed} & Ratio & \shortstack{Old/new\\status} \\
\midrule
\endfirsthead
\toprule
Input & $n$ & \shortstack{Old\\complete} & \shortstack{New\\complete} & \shortstack{Old\\observed} & \shortstack{New\\observed} & Ratio & \shortstack{Old/new\\status} \\
\midrule
\endhead
\midrule
\multicolumn{8}{r}{\footnotesize Continued on the next page.}\\
\endfoot
\bottomrule
\multicolumn{8}{@{}p{\linewidth}@{}}{\footnotesize All times are milliseconds. Completion medians and paired ratios are copied from the recorded JSON summaries. An unresolved arm has -- for completion time and no completion ratio; only then is its observed elapsed median shown. U = \texttt{UNKNOT}, K = \texttt{KNOTTED}, I = \texttt{INCONCLUSIVE}, ? = \texttt{UNKNOWN}. All measured outcomes for each arm agree within every row in this data set.}\\
\endlastfoot
\texttt{circle-\allowbreak 15} & 15 & 3.751 & 1.325 & -- & -- & 2.814 & U/U \\
\texttt{circle-\allowbreak 31} & 31 & 14.620 & 2.712 & -- & -- & 5.311 & U/U \\
\texttt{circle-\allowbreak 63} & 63 & 57.965 & 5.498 & -- & -- & 10.290 & U/U \\
\texttt{circle-\allowbreak 127} & 127 & 234.973 & 11.059 & -- & -- & 20.393 & U/U \\
\texttt{circle-\allowbreak 255} & 255 & 1\,036.076 & 22.204 & -- & -- & 46.720 & U/U \\
\texttt{circle-\allowbreak 511} & 511 & -- & 44.922 & 742.769 & -- & -- & I/U \\
\texttt{survivor-\allowbreak 00} & 16 & 7.648 & 7.130 & -- & -- & 0.964 & U/U \\
\texttt{survivor-\allowbreak 01} & 15 & 5.662 & 5.912 & -- & -- & 0.949 & U/U \\
\texttt{survivor-\allowbreak 02} & 15 & 6.428 & 6.446 & -- & -- & 1.009 & U/U \\
\texttt{survivor-\allowbreak 03} & 13 & 4.507 & 4.494 & -- & -- & 0.992 & U/U \\
\texttt{survivor-\allowbreak 04} & 13 & 3.771 & 3.755 & -- & -- & 0.996 & U/U \\
\texttt{survivor-\allowbreak 05} & 13 & 8.321 & 8.387 & -- & -- & 1.018 & U/U \\
\texttt{survivor-\allowbreak 06} & 12 & 2.876 & 2.971 & -- & -- & 0.974 & U/U \\
\texttt{survivor-\allowbreak 07} & 12 & 8.693 & 8.791 & -- & -- & 0.985 & U/U \\
\texttt{survivor-\allowbreak 08} & 11 & 2.977 & 2.947 & -- & -- & 1.004 & U/U \\
\texttt{survivor-\allowbreak 09} & 11 & 3.055 & 3.015 & -- & -- & 0.985 & U/U \\
\texttt{survivor-\allowbreak 10} & 11 & 2.807 & 2.597 & -- & -- & 1.083 & U/U \\
\texttt{survivor-\allowbreak 11} & 11 & 2.668 & 2.685 & -- & -- & 0.981 & U/U \\
\texttt{trefoil} & 3 & -- & -- & 0.712 & 0.625 & -- & I/I \\
\texttt{conway} & 11 & -- & -- & 7.777 & 7.630 & -- & I/I \\
\texttt{kinoshita\_\allowbreak terasaka} & 11 & -- & -- & 12.983 & 12.896 & -- & I/I \\
\texttt{hard\_\allowbreak unknot\_\allowbreak 8} & 8 & 1.277 & 1.440 & -- & -- & 0.891 & U/U \\
\texttt{grid\_\allowbreak determinant\_\allowbreak one\_\allowbreak knot} & 10 & -- & -- & 2.450 & 2.493 & -- & I/I \\
\texttt{unknot\_\allowbreak braid40} & 40 & 22.536 & 3.387 & -- & -- & 6.625 & U/U \\
\texttt{stress\_\allowbreak braid5\_\allowbreak 36} & 36 & -- & -- & 68.986 & 76.205 & -- & I/I \\
\texttt{monster} & 10 & 2.095 & 2.202 & -- & -- & 0.912 & U/U \\
\texttt{gst} & 48 & -- & -- & 789.078 & 756.285 & -- & I/I \\
\texttt{gordian} & 141 & 393.890 & 380.828 & -- & -- & 1.010 & U/U \\
\end{longtable}
\endgroup


\subsection{Full recognizer with ordinary filters and fixed allowances}

% Generated by generate_artifacts.py. This is a longtable, not a float.
\begingroup
\small
\setlength{\tabcolsep}{3pt}
\setlength{\LTcapwidth}{\linewidth}
\renewcommand{\arraystretch}{1.08}
\begin{longtable}{@{}>{\raggedright\arraybackslash}p{3.45cm}rrrrrrc@{}}
\caption{All full-recognizer benchmark inputs.}\label{tab:raw-pipeline-full}\\
\toprule
Input & $n$ & \shortstack{Old\\complete} & \shortstack{New\\complete} & \shortstack{Old\\observed} & \shortstack{New\\observed} & Ratio & \shortstack{Old/new\\status} \\
\midrule
\endfirsthead
\toprule
Input & $n$ & \shortstack{Old\\complete} & \shortstack{New\\complete} & \shortstack{Old\\observed} & \shortstack{New\\observed} & Ratio & \shortstack{Old/new\\status} \\
\midrule
\endhead
\midrule
\multicolumn{8}{r}{\footnotesize Continued on the next page.}\\
\endfoot
\bottomrule
\multicolumn{8}{@{}p{\linewidth}@{}}{\footnotesize All times are milliseconds. Completion medians and paired ratios are copied from the recorded JSON summaries. An unresolved arm has -- for completion time and no completion ratio; only then is its observed elapsed median shown. U = \texttt{UNKNOT}, K = \texttt{KNOTTED}, I = \texttt{INCONCLUSIVE}, ? = \texttt{UNKNOWN}. All measured outcomes for each arm agree within every row in this data set.}\\
\endlastfoot
\texttt{survivor-\allowbreak 00} & 16 & 9.901 & 10.219 & -- & -- & 0.997 & U/U \\
\texttt{survivor-\allowbreak 01} & 15 & 7.899 & 8.089 & -- & -- & 0.980 & U/U \\
\texttt{survivor-\allowbreak 02} & 15 & 8.498 & 8.977 & -- & -- & 0.968 & U/U \\
\texttt{survivor-\allowbreak 03} & 13 & 6.067 & 6.183 & -- & -- & 0.980 & U/U \\
\texttt{survivor-\allowbreak 04} & 13 & 5.256 & 5.425 & -- & -- & 0.962 & U/U \\
\texttt{survivor-\allowbreak 05} & 13 & 10.436 & 10.756 & -- & -- & 0.971 & U/U \\
\texttt{survivor-\allowbreak 06} & 12 & 4.238 & 4.314 & -- & -- & 0.982 & U/U \\
\texttt{survivor-\allowbreak 07} & 12 & 11.189 & 11.394 & -- & -- & 0.982 & U/U \\
\texttt{survivor-\allowbreak 08} & 11 & 4.159 & 4.241 & -- & -- & 0.974 & U/U \\
\texttt{survivor-\allowbreak 09} & 11 & 4.279 & 4.423 & -- & -- & 0.979 & U/U \\
\texttt{survivor-\allowbreak 10} & 11 & 3.971 & 4.051 & -- & -- & 0.980 & U/U \\
\texttt{survivor-\allowbreak 11} & 11 & 3.960 & 3.960 & -- & -- & 1.001 & U/U \\
\texttt{trefoil} & 3 & 0.036 & 0.034 & -- & -- & 1.053 & K/K \\
\texttt{conway} & 11 & 0.645 & 0.645 & -- & -- & 0.960 & K/K \\
\texttt{kinoshita\_\allowbreak terasaka} & 11 & 0.599 & 0.631 & -- & -- & 1.027 & K/K \\
\texttt{hard\_\allowbreak unknot\_\allowbreak 8} & 8 & 1.427 & 1.496 & -- & -- & 0.954 & U/U \\
\texttt{grid\_\allowbreak determinant\_\allowbreak one\_\allowbreak knot} & 10 & 0.069 & 0.064 & -- & -- & 1.005 & K/K \\
\texttt{unknot\_\allowbreak braid40} & 40 & 0.213 & 0.211 & -- & -- & 1.011 & U/U \\
\texttt{stress\_\allowbreak braid5\_\allowbreak 36} & 36 & 0.628 & 0.630 & -- & -- & 0.973 & K/K \\
\texttt{monster} & 10 & 2.993 & 3.297 & -- & -- & 0.903 & U/U \\
\texttt{gst} & 48 & 1.336 & 1.312 & -- & -- & 0.989 & K/K \\
\texttt{gordian} & 141 & -- & -- & 1\,015.181 & 1\,020.367 & -- & ?/? \\
\end{longtable}
\endgroup

